\section{Regularity required by the Gaussian remainder}
\label{sec:regularity}

Source Eq.~(35), p.~6, assumes a Gaussian packet and an
observable $f$ twice continuously differentiable near its mean, but
assigns the expansion a remainder $O(\norm{\Sigma}^2)$:
\begin{equation}
 \E{f(X)}=f(\mu)+\frac12\operatorname{tr}
              (\Sigma D^2f(\mu))+O(\norm{\Sigma}^2).
 \label{eq:source-observable-remainder}
\end{equation}
The stated differentiability does not imply that order, even for an
ideal Gaussian.

\begin{proposition}[A $C^2$ counterexample]
\label{prop:c2-remainder}
Set $\mu=1/2$ and $f(x)=|x-\mu|^3$ on $[0,1]$.
This observable is $C^2$, satisfies $0\leq f\leq1/8$, and has
$f(\mu)=f'(\mu)=f''(\mu)=0$. For a Gaussian of width $\sigma$ centered
at $\mu$, its expectation has leading order $\sigma^3$, rather than
$O(\sigma^4)$.
\end{proposition}

\begin{proof}
For an untruncated standard normal $Z$,
\begin{equation}
 \E{|\sigma Z|^3}=2\sqrt{\frac2\pi}\,\sigma^3.
 \label{eq:absolute-third-moment}
\end{equation}
The observable may be extended to a bounded $C^2$ function on $\R$
by a smooth cutoff that is one near $\mu$; this changes the displayed
expectation only by exponentially small terms as $\sigma\to0$.
Alternatively, normalize the Gaussian on $[0,1]$ itself. Its mean remains
$\mu$, its variance $\tau_\sigma^2$ satisfies
$\tau_\sigma^2\sim\sigma^2$, and the expectation in
\eqref{eq:absolute-third-moment} has only exponentially small changes.
Consequently
\begin{equation}
 \frac{\E{f(X)}}{\tau_\sigma^4}
 \sim 2\sqrt{\frac2\pi}\,\frac1\sigma\longrightarrow\infty.
 \label{eq:c2-ratio}
\end{equation}
The purported second-order covariance remainder is therefore not a
consequence of the stated $C^2$ hypothesis.
\end{proof}

The untruncated calculation is an exact Gaussian test; the compact
version addresses the localized packets used on a bounded computational
domain. Source Appendix~D.1 itself extends a narrow Gaussian integral
to the real axis. The exponentially small boundary correction cannot
repair the algebraic order in \eqref{eq:c2-ratio}.

A standard sufficient stronger hypothesis makes the intended order
precise. Here the operator norm of a fourth derivative means the norm
of its associated four-linear form.

\begin{proposition}[A sufficient fourth-derivative bound]
\label{prop:c4-remainder}
Let $X=\mu+Y$ with $Y\sim N(0,\Sigma)$ in fixed dimension $d$.
Suppose $f\in C^4(\R^d)$ and
$\sup_y\norm{D^4f(y)}_{\rm op}\leq M_4$. Then
\begin{align}
 \left|\E{f(X)}-f(\mu)-\frac12\operatorname{tr}
                         (\Sigma D^2f(\mu))\right|
 &\leq\frac{M_4}{24}
       \left[(\operatorname{tr}\Sigma)^2+
                         2\operatorname{tr}(\Sigma^2)\right]
       \label{eq:c4-remainder}\\
 &\leq\frac{M_4(d+2)}{24}\norm{\Sigma}_F^2.
       \nonumber
\end{align}
\end{proposition}

\begin{proof}
Taylor expansion through cubic order has pointwise remainder bounded
by $M_4\norm{Y}_2^4/24$. Centered Gaussian odd moments vanish, and
$\E{Y^\top D^2f(\mu)Y}=\operatorname{tr}(\Sigma D^2f(\mu))$.
The Gaussian fourth-moment identity is
$\E{\norm{Y}_2^4}=(\operatorname{tr}\Sigma)^2+
2\operatorname{tr}(\Sigma^2)$. Finally,
$(\operatorname{tr}\Sigma)^2\leq d\norm{\Sigma}_F^2$.
\end{proof}

Local $C^4$ hypotheses can be used with an explicit tail estimate and
control outside the local neighborhood; they cannot silently replace
the global bound in this proposition. This is a standard Taylor and
Gaussian-moment estimate, included to supply a repair rather than a
novel general approximation theorem. For the particular coordinate
observables $f=x_i$ and $f=x_ix_j$, the relevant moment identities are
exact for every distribution with the required moments. The regularity
counterexample therefore does not invalidate direct measurement of
means and covariance entries.
